\subsection{滤子基的聚点}

定义 2. 在拓扑空间 X 中，若点 x 属于 B 的所有集合的闭包，则称 x 为 X 上滤子基 B 的聚点。

若 \(x\) 是滤子基 \(\mathfrak{B}\) 的聚点，则由 § 6 第 3 目命题 4 的推论，它也是每个等价滤子基的聚点；特别地，\(x\) 是以 \(\mathfrak{B}\) 为基的滤子的聚点。

命题 3. 点 x 是滤子基 B 的聚点，当且仅当 x 的基本邻域系中的每个集合与 B 的每个集合相交。

这由定义立得。

这个命题与 § 6 第 2 目命题 1 的推论 2 表明，性质“x 是滤子 F 的聚点”等价于性质“存在一个滤子，它既细于 F 又细于 x 的邻域滤子”。换言之：

命题 4. 点 \(x\) 是滤子 \(\mathfrak{F}\) 的聚点，当且仅当存在细于 \(\mathfrak{F}\) 的滤子收敛于 \(x\)。

特别地，滤子 \(\mathfrak{F}\) 的每个极限点都是 \(\mathfrak{F}\) 的聚点。

推论. 超滤子 U 收敛于点 x，当且仅当 x 是 U 的聚点。

若 x 是滤子 F 的聚点，则它也是粗于 F 的每个滤子的聚点；同样，若把 X 的拓扑换成较粗的拓扑，则在新拓扑下 x 仍是 F 的聚点。

按定义，滤子基 \(\mathfrak{B}\)（\(\mathbf{X}\) 上的）的聚点所成之集就是集合 \(\bigcap_{\mathbf{M} \in \mathfrak{B}} \overline{\mathbf{M}}\)，由此得到

命题 5. 拓扑空间 X 上滤子基的聚点所成之集在 X 中是闭集。

命题 6. 设 \(\mathfrak{B}\) 是拓扑空间 X 的子集 A 上的滤子基。则 \(\mathfrak{B}\) 在 X 中的每个聚点都属于 \(\overline{\mathbf{A}}\)；反之，\(\overline{\mathbf{A}}\) 的每个点都是 A 上某个滤子的极限点。

第一个断言是显然的；另一方面，若 \(x \in \overline{A}\)，则 x 在 X 中的邻域滤子在 A 上的迹是 A 上的滤子，它显然收敛于 x。

注. 拓扑空间上的滤子可以没有聚点（更不用说极限点）；例如，在无限离散空间中，有限子集的余集所组成的滤子没有聚点。每个滤子都有聚点的空间在数学中起重要作用，我们将在 § 9 中研究它们。

